\section{Implementasi Spesifikasi Wajib}\label{sec:Implementasi Spesifikasi Wajib}

\subsection{Parser dan Validasi Input}\label{sub:Parser dan Validasi Input}

Parser berada pada package \texttt{parser}. Fungsi utama parser adalah
\texttt{ParseFile}. Fungsi ini menerima path berkas \texttt{.txt},
membaca seluruh isi file, lalu membentuk \texttt{BoardConfig}. Validasi
yang dilakukan meliputi:
\begin{itemize}[label=\textbullet]
    \item ekstensi file harus \texttt{.txt};
    \item baris pertama harus berisi tepat dua bilangan bulat positif
        \texttt{N M};
    \item jumlah baris grid harus sesuai dengan \texttt{N};
    \item setiap baris grid harus memiliki panjang \texttt{M};
    \item papan harus memiliki tepat satu start \texttt{Z} dan satu
        goal \texttt{O};
    \item checkpoint tidak boleh duplikat dan harus berurutan mulai
        dari \texttt{0};
    \item jumlah nilai cost harus tepat \texttt{N * M};
    \item setiap nilai cost harus berupa bilangan bulat tidak negatif.
\end{itemize}

Dengan validasi tersebut, solver dapat berasumsi bahwa struktur papan
yang diterima sudah konsisten.

\subsection{Pembangkitan Successor}\label{sub:Pembangkitan Successor}

Package \texttt{game} berisi aturan permainan. Fungsi
\texttt{GetSuccessors} membangkitkan semua langkah valid dari suatu
state dengan mencoba empat arah: atas, bawah, kiri, dan kanan.

\begin{lstlisting}[language=Go, basicstyle=\ttfamily\small, breaklines=true]
func GetSuccessors(board *parser.BoardConfig, s Node) []Edge {
    var successors []Edge
    for _, dir := range AllDirs {
        newState, cost, path, valid := Slide(board, s, dir)
        if valid {
            successors = append(successors, Edge{
                State:     newState,
                Cost:      cost,
                Direction: dir,
                Path:      path,
            })
        }
    }
    return successors
}
\end{lstlisting}

Fungsi \texttt{Slide} merupakan inti aturan permainan. Selama pemain
belum bertemu dinding, fungsi ini terus memindahkan posisi ke petak
berikutnya, menjumlahkan biaya, dan mencatat path petak yang dilewati.
Jika gerakan keluar papan atau melewati lava, gerakan langsung
ditolak. Jika bertemu checkpoint, checkpoint hanya diterima apabila
checkpoint tersebut merupakan target berikutnya atau checkpoint yang
sudah dilewati sebelumnya.

\subsection{Kerangka Pencarian}\label{sub:Kerangka Pencarian}

Solver menggunakan struktur \texttt{SearchNode} untuk menyimpan state
permainan dan informasi pencarian.

\begin{lstlisting}[language=Go, basicstyle=\ttfamily\small, breaklines=true]
type SearchNode struct {
    Node     graph.Node
    Parent   *SearchNode
    Move     graph.Dir
    GCost    int
    HCost    int
    Priority int
    StepPath []graph.Pos
    Index    int
}
\end{lstlisting}

Field \texttt{Parent}, \texttt{Move}, dan \texttt{StepPath} digunakan
untuk merekonstruksi solusi setelah goal ditemukan. Field
\texttt{GCost}, \texttt{HCost}, dan \texttt{Priority} digunakan oleh
priority queue.

Secara garis besar, algoritma pencarian berjalan sebagai berikut.
\begin{enumerate}
    \item Buat state awal dari posisi \texttt{Z} dengan
        \texttt{NextTarget = 0}.
    \item Masukkan state awal ke priority queue.
    \item Selama priority queue tidak kosong, ambil simpul dengan
        prioritas terkecil.
    \item Abaikan simpul jika sudah ada rute yang lebih murah menuju
        state yang sama.
    \item Jika simpul adalah goal, rekonstruksi path dan kembalikan
        hasil.
    \item Jika belum goal, bangkitkan semua successor valid.
    \item Untuk setiap successor, hitung biaya baru, heuristic, dan
        prioritas sesuai algoritma yang dipilih.
    \item Masukkan successor ke priority queue apabila rute tersebut
        lebih baik daripada rute sebelumnya.
\end{enumerate}

\subsection{Implementasi UCS, GBFS, dan A*}\label{sub:Implementasi Algoritma}

Tiga algoritma wajib diimplementasikan sebagai wrapper terhadap
\texttt{graphSearch}. Wrapper menjaga antarmuka pemanggilan tetap
sederhana.

\begin{lstlisting}[language=Go, basicstyle=\ttfamily\small, breaklines=true]
func UCS(board *parser.BoardConfig) Result {
    return graphSearch(board, AlgorithmUCS, "")
}

func GBFS(board *parser.BoardConfig, heuristicMode string) Result {
    return graphSearch(board, AlgorithmGBFS, heuristicMode)
}

func AStar(board *parser.BoardConfig, heuristicMode string) Result {
    return graphSearch(board, AlgorithmAStar, heuristicMode)
}
\end{lstlisting}

Perhitungan prioritas dilakukan pada satu fungsi agar perilaku ketiga
algoritma mudah dibandingkan.

\begin{lstlisting}[language=Go, basicstyle=\ttfamily\small, breaklines=true]
func calculatePriority(algorithm Algorithm, gCost, hCost int) int {
    switch algorithm {
    case AlgorithmUCS:
        return gCost
    case AlgorithmGBFS:
        return hCost
    case AlgorithmAStar:
        return gCost + hCost
    default:
        return gCost
    }
}
\end{lstlisting}

\subsection{Implementasi BFS dan DFS}\label{sub:Implementasi BFS DFS}

BFS dan DFS diimplementasikan sebagai algoritma bonus pada fungsi
\texttt{unweightedGraphSearch}. Fungsi ini menggunakan representasi
state dan successor yang sama dengan algoritma wajib, tetapi tidak
menggunakan priority queue. BFS mengambil elemen dari depan slice
seperti queue, sedangkan DFS mengambil elemen terakhir seperti stack.

\begin{lstlisting}[language=Go, basicstyle=\ttfamily\small, breaklines=true]
func BFS(board *parser.BoardConfig) Result {
    return unweightedGraphSearch(board, AlgorithmBFS, nil)
}

func DFS(board *parser.BoardConfig) Result {
    return unweightedGraphSearch(board, AlgorithmDFS, nil)
}
\end{lstlisting}

Pada BFS, kedalaman simpul menyatakan jumlah slide yang telah
dilakukan. Oleh karena itu, BFS optimal terhadap jumlah slide. Namun,
karena cost setiap slide bergantung pada total cost tile yang dilewati,
BFS tidak selalu menghasilkan total cost terkecil. DFS digunakan
sebagai pembanding karena dapat menemukan solusi dengan strategi
penelusuran yang berbeda, tetapi tidak menjamin optimalitas.

\subsection{Heuristic}\label{sub:Implementasi Heuristic}

Heuristic digunakan oleh GBFS dan A*. Program menyediakan tiga pilihan.
\begin{enumerate}
    \item \texttt{H1}: Manhattan normal ke target berikutnya.
    \item \texttt{H2}: alternatif heuristic Manhattan ke target
        berikutnya dikali biaya minimum tile yang dapat dilewati.
    \item \texttt{H3}: alternatif heuristic total Manhattan dari
        posisi saat ini menuju semua checkpoint tersisa secara
        berurutan, lalu menuju goal.
\end{enumerate}

\begin{lstlisting}[language=Go, basicstyle=\ttfamily\small, breaklines=true]
func Heuristic(board *parser.BoardConfig, node graph.Node, mode string) int {
    switch mode {
    case HeuristicManhattan:
        return manhattanToNextTarget(board, node)
    case HeuristicManhattanCost:
        return manhattanToNextTarget(board, node) * minTraversableCost(board)
    case HeuristicRemaining:
        return remainingManhattan(board, node)
    default:
        return manhattanToNextTarget(board, node) * minTraversableCost(board)
    }
}
\end{lstlisting}

\subsection{Rekonstruksi Solusi}\label{sub:Rekonstruksi Solusi}

Ketika goal ditemukan, solver tidak langsung memiliki urutan gerakan
dari awal ke akhir. Setiap \texttt{SearchNode} hanya menyimpan parent.
Oleh karena itu, program melakukan traversal mundur dari goal ke start,
menyimpan simpul dan gerakan yang dilewati, lalu membalik urutan hasil.

Hasil akhir disimpan dalam struktur \texttt{Result}.
\begin{lstlisting}[language=Go, basicstyle=\ttfamily\small, breaklines=true]
type Result struct {
    Found       bool
    Moves       []graph.Dir
    PathNodes   []graph.Node
    StepPaths   [][]graph.Pos
    TotalCost   int
    Iterations  int
    ExecutionMs int64
}
\end{lstlisting}

\subsection{CLI}\label{sub:CLI}

CLI berada pada \texttt{src/main.go}. Program meminta pengguna
memasukkan path file testcase, pilihan algoritma, dan pilihan
heuristic jika algoritma yang dipilih adalah GBFS atau A*. Setelah
solver selesai, CLI menampilkan status solusi, urutan gerakan, total
cost, jumlah iterasi, waktu eksekusi, dan visualisasi papan untuk
setiap langkah solusi.

\subsection{Penyimpanan Solusi}\label{sub:Penyimpanan Solusi}

Setelah hasil pencarian ditampilkan, CLI menanyakan apakah pengguna
ingin menyimpan hasil ke berkas \texttt{.txt}. Jika pengguna memilih
ya, program meminta path output. Ekstensi \texttt{.txt} ditambahkan
secara otomatis apabila pengguna belum menuliskannya.

Berkas solusi yang disimpan berisi ringkasan algoritma, heuristic
jika digunakan, status solusi, urutan gerakan, total cost, jumlah
iterasi, waktu eksekusi, dan visualisasi papan untuk setiap langkah.
Jika solusi tidak ditemukan, program tetap dapat menyimpan status
\textit{not found}, jumlah iterasi, waktu eksekusi, dan papan awal.
